\section{Conclusion}
\label{sec:conclusion}

We presented the first systematic six-method comparison on three
Ekman-mapped emotion corpora (GoEmotions, ISEAR, WASSA-21) under two
evaluation protocols (mixed and leave-one-dataset-out), with three to
five seeds per cell and paired bootstrap significance testing.

Our headline finding is counter-intuitive: source-only training with
focal loss replacing cross-entropy produces the largest measured
cross-dataset gain (\focalwassagain{} F1 points on WASSA-21 LODO,
$p = \focalwassap$), outperforming both domain-adversarial methods
and their focal-loss combinations. Two of three LODO targets are
method-invariant — no method improves on the CE baseline within seed
variance. The pattern is consistent with class-distribution shift
being the dominant cross-dataset gap in emotion classification, which
focal loss directly targets but adversarial feature alignment does not.

Negative findings are characterized rather than hidden: DANN training
is unstable at moderate $\lambda$, CDAN produces stable but
non-generalizing source-source alignment, and adversarial-focal
combinations underperform focal alone. A DeBERTa-v3-large robustness
validation preserves the qualitative finding under elevated seed
variance that we attribute to the small WASSA-21 test set.

The practical recommendation we draw is concrete: for cross-corpus
text classification problems with severe class imbalance, the simplest
class-imbalance-aware loss is a stronger baseline than off-the-shelf
adversarial domain adaptation. Future work should examine whether
this conclusion extends to other cross-dataset NLP benchmarks with
similar imbalance properties, and whether a method that explicitly
combines focal-style class-imbalance handling with target-aware
(rather than source-source) feature alignment can produce gains
beyond focal alone.
